\honest{Humans in Funny Suits}{Eliezer Yudkowsky}{2008}

On screen, the human species travels to the stars and finds aliens who look ``remarkably like humans in funny suits.'' I pretend to explain this by convergent evolution, never mind that a kangaroo resembles us less than a chimp, a frog less than a kangaroo, and insects, with six legs, outside skeletons and different eyes, less still; a real alien might not even use DNA or proteins. Any intelligence needs hands, so two legs and an upright walk; binocular vision, so two eyes in a head; speech, so ears and lips; a nose, because a face without one would look wrong to mates; attractive females, since ugly aliens would breed less; and English, or they could not build a civilization. Then the real explanation: costumed humans are ``the easy way out.''

But the real problem is mind. An alien is a human in a funny suit if it thinks like a human, especially an English speaker of our time, even if it is an angular creature of pure crystal. The Americans in \textsc{Psycho} (1960) seemed to me more alien than most screen aliens. To write a culture unlike your own you must see yours as a special case, not the norm. History helps, but it is only black letters on white pages. Living a year in China or Dubai or among the !Kung would help more, I suspect, though ``this I have never done, being busy.'' Seeing your humanity as a special case is much harder.

In every known culture humans ``seem to experience'' the same six emotions, and these are shown ``by the same facial expressions.'' \nb{The second claim was disputed well before the post. A 1994 review by James Russell argued that the supporting studies' methods may have shaped their results, and a 2019 review by Barrett and colleagues found that expressions vary substantially across cultures.} The next time an alien, or an AI, on screen gets angry, and it will, I bet it shows the human facial expression for anger. Humans are alike because, in a sexual species, complex adaptations must be shared or they would not assemble. (Do the aliens reproduce sexually? Swap genes like bacteria? Form colonies like fungi?)

Our ancestors had to model only minds like their own, so ``we evolved to predict Other Minds by putting ourselves in their shoes.'' You may protest that you do not assume others are like you; they are angry when you are sad, and believe other things. But a brain is too complex to model ``neuron-by-neuron.'' A system that complex and unlike you would take scientific lifetimes to unravel, and you cannot build a model that predicts others as well as you do. The only way to grasp another is to run a shadow of its anger and beliefs on your own: if you were angry at you, and believed you godless scum, you would try to hurt you. I call this ``empathic inference.'' \nb{This is simulation theory, one side of a debate. Its rival holds that people use a tacit theory of mind, an abstract model of goals and beliefs, which the post's argument does not consider. ``Belief in Intelligence,'' the previous post in the book, uses such a model.}

Empathic inference fails for minds unlike ours. I can tell you to imagine an alien raised in four spatial dimensions, but you cannot rewire your visual cortex to see as it would, and neither of us can feel an alien's emotions. An alien with no sense of humor would not see the point of the Marx Brothers; ``you've never antled.'' Readers will say evolution must give every intelligence a sense of humor, or that my jokes are not funny enough; I compare the second to shouting English abroad. \nb{The first objection gets no answer.} Laughter may be valuable even if not all possible minds share it, our own special part of the gift we give to tomorrow. Universalizability among all possible minds is one idea about ethics I cannot rescue.

The artists who drew aliens carrying off girls in torn dresses ``did not make that error by reasoning''; it ``just seemed to them'' that the girl was sexy, as a property of the girl and the dress. Your words have meaning and your jokes are funny. What has that to do with aliens? We do not think of ourselves as human while being human, and no act of will removes it; even imagining a slime monster that mates with slime monsters, it is hard to imagine it not preferring the girl. Do not explain why any intelligent mind must have emotions. Natural selection builds complex machines without them. It is a real alien, an optimizer that does not work as we do. Much of biology's progress since the 1960s has been enforcing a ban on anthropomorphizing evolution, won, I suspect, only by crushing experiments and clear mathematics; I have fought the same battle about AI for years. \nb{Yudkowsky calls evolution ``stupid in an absolute sense'' in ``Thou Art Godshatter,'' and does not count it as an intelligence. It shows that making complex machines does not need emotions, not that intelligence does not.}

Writing real aliens, not merely conveniently incomprehensible ones, is the proverbial test of a science fiction author. Jack Vance's humans from other cultures are more alien than most aliens; start with \textsc{City of the Chasch}, and \textsc{The Mote in God's Eye} also counts. The low point, which I think Orson Scott Card named, was the Star Trek episode whose aliens had rewritten the preamble to the U.S. Constitution. This Great Failure of Imagination is not only about science fiction or AI. ``The inability to imagine the alien is the inability to see yourself.'' Who can see a human camouflaged against a human background?
